
\subsubsection{Dataset}

We constructed a 150-model heterogeneous zoo from publicly available pretrained models:

\begin{itemize}
\item \textbf{CNNs} (50 models): torchvision pretrained models (ResNet, EfficientNet, DenseNet variants) originally trained on ImageNet
\item \textbf{Vision Transformers} (50 models): timm library ViT models originally pretrained on ImageNet
\item \textbf{MLPs} (50 models): Fully-connected networks trained on MNIST
\end{itemize}

Generalization gaps were computed as (test accuracy - train accuracy) evaluated on CIFAR-10. ImageNet evaluation was not performed due to dataset unavailability in the experimental environment. This creates domain shift (ImageNet pretraining evaluated on CIFAR-10 transfer), which may affect measured generalization characteristics.

Train/validation/test split: 150/30/30 models, stratified by architecture family. Generalization gap range: [-0.12, 0.31]. Average parameters per model: ~25M (CNNs), ~86M (ViTs), ~2M (MLPs).

The 150-model scale represents a proof-of-concept experiment, reduced from an initially planned 750-model full-scale validation (600 train, 150 test).

\subsubsection{Baselines}

\textbf{Flat-Weight MLP}: Concatenates all model weights into a single vector with zero-padding for variable-length inputs. Processes through a 3-layer MLP with the same training protocol as CAWE. This tests whether architecture-specific tokenization provides value beyond naive vectorization.

\textbf{Random Forest with Engineered Features}: Extracts hand-crafted weight statistics (layer-wise L2 norms, sparsity ratios, spectral radius of weight matrices) and trains a sklearn RandomForestRegressor with 100 trees. This tests whether learned representations outperform explicit feature engineering.

\subsubsection{Validation Components}

Five validation components with prespecified success thresholds:

\textbf{Component 1 (Per-Family Ablation)}: Train CAWE separately on CNN-only, Transformer-only, and MLP-only subsets. Success: $\rho$ > 0.7 for all three families. This tests whether tokenization preserves family-specific quality signals.

\textbf{Component 2 (Architecture Clustering)}: Extract CAWE embeddings for all test models and compute silhouette score with architecture family labels. Success: silhouette > 0.5. This tests whether shared NFT maintains architecture-aware representations.

\textbf{Component 3 (Flat-Weight Baseline)}: Compare CAWE vs flat-weight MLP using paired t-test. Success: $\Delta\rho$ > 0.15 with p < 0.001.

\textbf{Component 4 (Random Forest Baseline)}: Compare CAWE vs random forest using Wilcoxon signed-rank test. Success: $\Delta\rho$ > 0.1 with p < 0.01.

\textbf{Component 5 (Robustness)}: Test CAWE with token dimensions D $\in$ {64, 128, 256, 512}. Success: $\geq 2$ variants achieve $\rho$ > 0.65.

Gate criterion: $\geq 3/5$ components pass thresholds.

\subsubsection{Implementation}

Framework: PyTorch 2.0. Hardware: Single NVIDIA GPU. NFT Library: nfn (PyPI). Training time: ~2 hours for 150-model experiment.

\subsubsection{Evaluation Metrics}

Primary metric: Spearman rank correlation $\rho$ with bootstrap 95\% confidence intervals (1000 bootstrap samples). Clustering metric: Silhouette score. Statistical tests: paired t-test (flat-weight baseline), Wilcoxon signed-rank test (random forest baseline).
